\editionPage[\pdfbookmark[0]{Indications des défis}{bm-indices-debut-parent}]{indices-debut}{Indications progressives}{Défis 1 à 3}{Ne lire que l’aide nécessaire.}
\begin{blocEdition}{indices-debut}
\editionEtape{D01 · Une fraction entre deux autres}\textbf{Indice 1.} Mettre les deux bornes au dénominateur 60.\par\medskip\textbf{Indice 2.} Chercher les numérateurs entiers strictement entre les deux bornes.\par
\editionEtape{D02 · Comparer sans grands calculs}\textbf{Indice 1.} Écrire $2^{30}$ comme un cube.\par\medskip\textbf{Indice 2.} Élever au carré deux nombres positifs conserve leur ordre.\par
\editionEtape{D03 · Une recherche entière exhaustive}\textbf{Indice 1.} Décomposer 180 en facteurs premiers.\par\medskip\textbf{Indice 2.} Dans un carré, les exposants premiers sont pairs.\par

\end{blocEdition}
\editionPage{indices-4}{Indications progressives}{Défis 4 à 6}{Ne lire que l’aide nécessaire.}
\begin{blocEdition}{indices-4}
\editionEtape{D04 · Un motif qui cache une identité}\textbf{Indice 1.} Développer les deux produits.\par\medskip\textbf{Indice 2.} Utiliser l’identité avec x = 1000.\par
\editionEtape{D05 · Deux solutions et un contexte}\textbf{Indice 1.} Calculer aire du carré moins aire du rectangle.\par\medskip\textbf{Indice 2.} Une différence de valeur absolue 1 correspond à deux égalités possibles.\par
\editionEtape{D06 · Promotions dans quel ordre ?}\textbf{Indice 1.} Écrire les deux prix avec P.\par\medskip\textbf{Indice 2.} Développer $0{,}8(P-10)$.\par

\end{blocEdition}
\editionPage{indices-7}{Indications progressives}{Défis 7 à 9}{Ne lire que l’aide nécessaire.}
\begin{blocEdition}{indices-7}
\editionEtape{D07 · Une rencontre de débits}\textbf{Indice 1.} Calculer le débit net pendant le remplissage.\par\medskip\textbf{Indice 2.} Distinguer la durée de vidange de la durée totale.\par
\editionEtape{D08 · Une fonction cachée derrière des essais}\textbf{Indice 1.} Ajouter une expression nulle en -1, 0 et 1.\par\medskip\textbf{Indice 2.} Le produit $(x+1)x(x-1)$ est nul pour les trois entrées.\par
\editionEtape{D09 · Trois tarifs, un choix argumenté}\textbf{Indice 1.} Comparer A et B, B et C, puis A et C.\par\medskip\textbf{Indice 2.} Tester les entiers proches des points d’égalité.\par

\end{blocEdition}
\editionPage{indices-10}{Indications progressives}{Défis 10 à 12}{Ne lire que l’aide nécessaire.}
\begin{blocEdition}{indices-10}
\editionEtape{D10 · Faire baisser une moyenne sans baisser la médiane}\textbf{Indice 1.} Fixer d’abord trois petites valeurs dont la troisième est 8.\par\medskip\textbf{Indice 2.} Baisser une valeur extrême de 5 sans lui faire franchir le rang central.\par
\editionEtape{D11 · Un jeu de somme équitable ?}\textbf{Indice 1.} Compter les couples pour chacune des trois sommes.\par\medskip\textbf{Indice 2.} Se limiter aux couples qui désignent un gagnant.\par
\editionEtape{D12 · Une corde dans un cercle}\textbf{Indice 1.} Appliquer Pythagore à OHA.\par\medskip\textbf{Indice 2.} La base AB a pour hauteur OH.\par

\end{blocEdition}
\editionPage{indices-13}{Indications progressives}{Défis 13 à 15}{Ne lire que l’aide nécessaire.}
\begin{blocEdition}{indices-13}
\editionEtape{D13 · Une droite parallèle partage une aire}\textbf{Indice 1.} Le rapport des aires est le carré du rapport des longueurs.\par\medskip\textbf{Indice 2.} Le rapport recherché est positif dans cette configuration.\par
\editionEtape{D14 · Deux angles de visée}\textbf{Indice 1.} Si la distance B-pied vaut x, la visée à 45 degrés donne une hauteur x.\par\medskip\textbf{Indice 2.} L’autre visée donne $x=(x+10)\tan30^\circ$.\par
\editionEtape{D15 · La même aire après deux changements}\textbf{Indice 1.} Suivre les facteurs d’aire, pas la position.\par\medskip\textbf{Indice 2.} L’aire finale vaut $60	imes4r^2$.\par

\end{blocEdition}
\editionPage{indices-16}{Indications progressives}{Défis 16 à 18}{Ne lire que l’aide nécessaire.}
\begin{blocEdition}{indices-16}
\editionEtape{D16 · Une section d’aire imposée}\textbf{Indice 1.} Retrouver le rayon de section grâce à son aire.\par\medskip\textbf{Indice 2.} Dans le triangle rectangle des rayons, utiliser $10^2=OH^2+8^2$.\par
\editionEtape{D17 · Un script de recherche et sa preuve}\textbf{Indice 1.} Calculer les sommes successives $1+2+\cdots+k$.\par\medskip\textbf{Indice 2.} La variable n est augmentée après son ajout à s.\par
\editionEtape{D18 · Le meilleur rectangle entier}\textbf{Indice 1.} La somme des deux côtés vaut 23.\par\medskip\textbf{Indice 2.} Comparer $x(23-x)$ à $132$ ou écrire $(x-11)(x-12)$.\par

\end{blocEdition}
